% Reviewed copy of FF.26, integrated in F_section.tex after two independent audits. Do not input this duplicate.
\begin{claim}{FF.26}{Proved}\label{clm:FFstrictreconstruct}
For every prime \(p\) and every finite extension
\(K/\operatorname{Frac}(\mathbb Z_{(p)}^{sh})\),
the field \(K\) satisfies \(\mathcal R(K)\).
Consequently every subfield of such a \(K\) reconstructs.
\end{claim}

\begin{proof}
Throughout, all topoi and their structural isomorphisms are
relative to the stated field. Finite scalar extension of a
pinned map is formed by the corresponding finite \'etale
slice. We may make finite extensions during the proof:
if a unit on \(X_L\) realizes the original point function
over a finite extension \(L/K\), pass to the finite normal
closure of \(L/K\). Every conjugate unit has the same
geometric point function. Equality on geometric points of
the reduced scheme makes the units equal, so the unit
descends to \(K\). This reduction does not assume
reconstruction over any of the fields in advance.

Let \(X/K\) be regular, quasiprojective and connected, with
\(x_0\in X(K)\), and let
\[
 \varphi:X_{\mathrm{\acute et}}\longrightarrow
                    \mathbb G_{m,K,\mathrm{\acute et}}
\]
be pinned. Regularity and the rational point make \(X\)
geometrically integral; in dimension zero it is \(\operatorname{Spec}K\),
where Claim~\ref{clm:FF1} applies directly. We henceforth
assume positive dimension. Its rational points are Zariski
 dense. Indeed, an \'etale coordinate chart at \(x_0\)
and Hensel's lemma lift every point of a sufficiently
small valuation polydisc in \(K^n\) to \(X(K)\).
The image of a proper closed subset under the quasi-finite
coordinate map has dimension less than \(n\). A product
of infinite valuation balls cannot lie in a proper
algebraic subset of \(\mathbb A^n_K\), by induction
on the variables of a nonzero polynomial. Thus every
nonempty open of \(X\) has a rational point. This
also holds after each finite extension and permits the
same argument on all inverse-image opens below. Write \(b(x)\) for its actual
value at a finite-extension point, and \(b_0=b(x_0)\).
Its pulled-back \(p\)-Kummer class is
\[
                   \xi=\varphi^*\kappa_p(T)
                          \in H^1(X,\mathbb Z_p(1)).
\]
In particular, evaluating \(\xi\) at \(x\in X(L)\)
gives \(\kappa_p(b(x))\), for every finite \(L/K\).

\emph{1. Integral boundary residues.}
We first prove a calculation requiring only a compatible
\(p\)-Kummer class with actual values at all finite-extension
points, rather than an entire topos map. On a punctured
rational line, after splitting its punctures over a finite
extension, normalization at a rational point writes the
class as
\[
 \kappa_p(c)+\sum_{i=1}^r a_i
       \kappa_p\left(\frac{T-r_i}{t_0-r_i}\right),
                  \qquad a_i\in\mathbb Z_p.             \tag{FF.26.1}
\]
This follows from the Kummer sequence and the vanishing of
the Picard group of a localization of \(K[T]\).
Let \(\pi\) be a uniformizer of this finite extension,
and evaluate at \(r_i+\pi^N\) and \(r_i+\pi^{N+1}\),
where \(N\) is larger than the valuations of all nonzero
root differences. Applying the valuation map on
\(H^1(K,\mathbb Z_p(1))\), the difference between these
two evaluations is \(a_i\): the other factors have
constant valuation. Both actual values have ordinary
integer valuation. Thus \(a_i\in\mathbb Z\). The residue
at infinity, the negative sum of the finite residues,
is integral as well.

For a smooth curve and a chosen rational boundary point
\(q\), Riemann--Roch gives a function whose only pole
is \(q\). Remove its branch values and the images of
all boundary points. The resulting map to a punctured
line is finite \'etale. Corestriction of the class has
actual values: at each point of the line these are the
products of the field norms of the actual values in its
fibre. The residue of a norm is the sum of the residue
degrees times the source residues. It has no factor
equal to the ramification index. Since \(q\) is the
only point above infinity, that residue is exactly
the residue at \(q\). Hence it is integral by (FF.26.1).
This argument is valid after making a boundary point
rational by a finite extension.

Choose a smooth projective compactification \(\bar X\)
with normal-crossing boundary, available in characteristic
zero. At a smooth point of a chosen boundary component,
away from the other components, take \'etale coordinates
in which that component is the first coordinate divisor.
Fix the other coordinates and take the curve component
through the point. Its local intersection multiplicity
with the chosen divisor is one. Normalize its closure
and apply the preceding curve argument. Thus each residue
\(m_i\) of \(\xi\) at a geometric boundary component
\(D_i\) belongs to \(\mathbb Z\). After another finite
extension all components and their residue data are defined
over the base. Put \(D=\sum_i m_iD_i\).

\emph{2. Forcing arithmetic equivariance over a finite local field.}
Let \(j:X\to A\) be the generalized Albanese map based
at \(x_0\), so \(j(x_0)=0\). The geometric part of
\(\xi-\kappa_p(b_0)\) is a continuous integral functional
\[
                  \lambda:T_pA\longrightarrow\mathbb Z_p(1)
\]
equivariant for \(H_0=G_K\). The Albanese comparison and
its finite obstruction are exactly those in
\citet[Proposition~4.1 and Corollary~4.2]{FF_SpiessSzamuely2003};
the finite obstruction disappears on taking the inverse
limit, since the Tate module of a finite group is zero.
The map \(j(X_{\bar K})\) generates \(A_{\bar K}\).

Choose a number field \(D_0\subset K\) containing all
the finite displayed geometric data and \(b_0\), and
such that \(D_0K_p=K\). Such a field exists because
\(K/K_p\) is finite. The chosen prolongation of the
\(p\)-adic valuation gives a finite extension
\(M_0=(D_0)_v\) of \(\mathbb Q_p\). The field \(K\)
is the fraction field of the corresponding strict
henselization of \(D_0\): each of these two strict
henselian fields contains the other by the universal
property. The completion/Krasner comparison therefore
identifies
\[
                  H_0=G_K=I_{M_0}\triangleleft G_{M_0}.
                                                               \tag{FF.26.2}
\]
These are compatible identifications for the Tate modules
of the varieties defined over \(D_0\); their torsion
points and all relevant point-Kummer cocycles are algebraic
over the number field.

Apply Claim~\ref{clm:FFpdetector} to the finite-dimensional
space
\[
                  S=\operatorname{Hom}_{H_0}
                       (V_pA,\mathbb Q_p(1)).
\]
Choose finitely many points \(x_i\), over one finite
extension \(E/K\), whose Kummer evaluations jointly
detect \(S\). Include their actual values \(b(x_i)\)
in the finite data, and choose a number field
\(D_1\subset E\) containing \(D_0\), all these data,
and a generator for \(E/K_p\). Then
\(E=D_1K_p\), and with \(M_1=(D_1)_v\) we have
\[
        H_1=G_E=I_{M_1}\triangleleft G_{M_1},
                  \qquad G_{M_1}\subset G_{M_0}.         \tag{FF.26.3}
\]
Restriction to \(H_1\) preserves joint detection,
by restriction--corestriction over \(\mathbb Q_p\).

For \(g\in G_{M_1}\), let
\(g\lambda=\chi_p(g)\lambda\rho_A(g)^{-1}\).
The normality in (FF.26.2) shows that
\(\delta_g=g\lambda-\lambda\) still belongs to \(S\).
Put
\[
 c_i=\kappa_A(j(x_i))\in H^1(G_{M_1},V_pA),\qquad
 d_i=\kappa_p(b(x_i)/b_0)\in H^1(G_{M_1},\mathbb Q_p(1)).
\]
The actual point-value equality is
\[
                \lambda_*\operatorname{res}_{H_1}(c_i)
                        =\operatorname{res}_{H_1}(d_i). \tag{FF.26.4}
\]
An \(H^1\)-class restricted from a group to a normal
subgroup is invariant under that group's conjugation
action, including its action on the coefficient module.
Apply this elementary cocycle identity to both \(c_i\)
and \(d_i\) in (FF.26.4). It gives
\[
               (\delta_g)_*\operatorname{res}_{H_1}(c_i)=0
                                      \quad\hbox{for all }i.
\]
Joint detection proves \(\delta_g=0\). Thus \(\lambda\)
is \(G_{M_1}\)-equivariant. This argument uses cohomology
classes, so the conjugation coboundaries cause no problem;
it does not require a trivial cyclotomic character and
does not assert that arbitrary topos maps descend to
\(M_1\).

\emph{3. A power of the class is the class of a unit.}
We give the integral Picard calculation to avoid an
incorrect local Tate-Hom assertion for general 1-motives.
Let \(P=\operatorname{Pic}^0_{\bar X}\), and let
\(B_X\subset\bigoplus_i\mathbb ZD_i\) be the group
of boundary divisors algebraically equivalent to zero.
The Picard 1-motive is \([B_X\to P]\), dual to
\([0\to A]\). Its realization gives
\[
 0\longrightarrow T_pP\longrightarrow
 H^1(X_{\bar K},\mathbb Z_p(1))\longrightarrow
 B_X\otimes\mathbb Z_p\longrightarrow0.                \tag{FF.26.5}
\]
Here the last map sends the geometric class to its
boundary residues. The construction, including the
finite obstruction before inverse limit, is given in
the proof of \citet[Proposition~4.1]{FF_SpiessSzamuely2003}.
The field has characteristic zero, so its prime-to-characteristic
statement includes the present prime \(p\).

By Step~1, the residue vector \(D\) is integral, while
(FF.26.5) puts it in \(B_X\otimes\mathbb Z_p\).
Consequently its class in the finitely generated group
\((\bigoplus_i\mathbb ZD_i)/B_X\) is torsion of order
prime to \(p\). Choose an integer \(N>0\), prime
to \(p\), such that \(ND\in B_X\).
Step~2 says that the geometric class is
\(G_{M_1}\)-invariant. The existence of its invariant
multiple lifting \(ND\) in (FF.26.5) means
\[
                 [ND]\in\bigcap_n p^nP(M_1).           \tag{FF.26.6}
\]
One can check this directly in the 1-motive, without
any Hom theorem. At level \(p^n\), a lift of \(ND\)
is represented by a \(p^n\)-division point of
\([ND]\in P\). The fixed residue vector allows us
to choose its lattice component equal to \(ND\).
Galois invariance then makes that division point
\(M_1\)-rational: a change of representative would
change the lattice component by \(p^n b\), and
\(p^n b=0\) implies \(b=0\) in the torsion-free
lattice. This holds for every \(n\).

For an abelian variety over a finite extension \(M_1\)
of \(\mathbb Q_p\), the compact group \(P(M_1)\)
contains an open subgroup isomorphic, under its formal
logarithm, to a finite free \(\mathbb Z_p\)-module.
Its infinitely \(p\)-divisible subgroup is therefore
exactly its finite prime-to-\(p\) torsion subgroup.
For completeness, modulo the open subgroup there is
only a finite group. A fixed power of \(p\) kills
its \(p\)-primary part; after multiplying by its
prime-to-\(p\) exponent, an infinitely \(p\)-divisible
element lies in every \(p^n\) multiple of the free
\(\mathbb Z_p\)-module and is zero. The converse for
prime-to-\(p\) torsion is immediate.
Thus (FF.26.6) is killed by an integer \(d_0>0\)
prime to \(p\), and \(dD\) is principal, where
\(d=Nd_0\).

All divisor and Picard data were defined over the
number field \(D_0\). Their equality after extension
to \(M_1\) implies the equality over \(D_0\), since
extension of scalars is injective on points of a scheme.
The principal divisor has a rational function over
\(D_0\), uniquely after fixing its value at \(x_0\):
the one-dimensional space of functions with that divisor
descends, or equivalently one applies Hilbert 90 to its
scalar ambiguity. Choose
\(h\in\mathcal O(X)^\times\) with
\(\operatorname{div}_{\bar X}(h)=dD\) and
\(h(x_0)=b_0^d\).

The difference between \(d\lambda\) and the geometric
Kummer class of \(h\) has zero residues. It belongs,
by (FF.26.5), to \((T_pP)^{G_{M_1}}\), which is zero:
the torsion subgroup of \(P(M_1)\) is finite, by
the same compact Lie-group argument. Hochschild--Serre
and evaluation at \(x_0\) then give the integral identity
\[
                         d\xi=\kappa_p(h).              \tag{FF.26.7}
\]
In particular, for every point over a finite extension,
\[
                         b(x)^d=\zeta_xh(x),
                         \qquad\zeta_x\in\mu_{p'}\subset K_p.
                                                               \tag{FF.26.8}
\]
This uses the integral kernel (FF.24.1), not merely
the kernel after tensoring cohomology with \(\mathbb Q_p\).

\emph{4. Curves with large-degree points and a resultant.}
We record explicitly two elementary tools. Every nonempty
open of a geometrically integral curve over a finite
extension of \(K_p\) has closed points of unbounded
degree. Choose a nonconstant map from its proper
normalization to \(\mathbb P^1\). The polynomials
\(T^n-\pi\) give closed points of degree \(n\) on
the line, by Eisenstein's criterion. For arbitrarily
large \(n\) their support avoids the finite image
of the omitted points; any point in the fibre has
degree at least \(n\).

Second, if \(V\subset K(X)\) is a finite-dimensional
space of rational functions regular on a nonempty open
\(U\), after a finite extension there is a smooth
geometrically integral curve mapping into \(U\) on
which restriction is injective on \(V\). Choose finitely
many geometric points of \(U\) whose evaluations are
injective on \(V\); this is possible by successively
testing a nonzero function in the remaining kernel.
Put them over a finite extension. A geometrically
irreducible curve through them exists by
\citet[Theorem~1.8, Corollary~1.9]{FF_CharlesPoonen2016},
whose hypotheses allow an arbitrary base field and
a finite set of closed points. Normalize its reduced
structure. In characteristic zero the result is a
smooth curve. Evaluation at its points above the
chosen set proves the asserted injectivity.

Suppose now that the geometric class \(\lambda\)
is nonzero. Then \(h\) in (FF.26.7) is nonconstant.
The actual point function is nonconstant as well:
if all values were equal to a base-field constant,
(FF.26.8) would put all values of \(h\) in the
base field. A nonconstant regular function cannot
have this property, by restricting to a curve and
using closed points of arbitrarily large degree
on its image.
Thus \(U=\varphi^{-1}(\mathbb G_m\setminus\{1\})\)
is a nonempty open of \(X\). Apply Steps~1--3 on
\(U\) to the pullback of \(T-1\), using a
rational point of that open, which exists by the preceding
density argument.
They give a positive integer \(e\), a rational
function \(g\) that is a unit on \(U\), and
\[
                     (b(x)-1)^e=\eta_xg(x),
                           \qquad\eta_x\in\mu_{p'}.
                                                               \tag{FF.26.9}
\]
Consider the finite-dimensional space generated by
\(h^ig^j\), for \(0\leq i\leq e\) and
\(0\leq j\leq d\). Choose a curve as in the preceding
paragraph, so restriction to that curve is injective
on this space. Write
\[
              B=e\deg\operatorname{div}_\infty(h|_C)
                   +d\deg\operatorname{div}_\infty(g|_C).
\]
Choose a closed point \(x\in C\) of degree greater
than \(B\). The rational function
\[
 R=\operatorname{Res}_Z
       \bigl(Z^d-\zeta_xh,\ (Z-1)^e-\eta_xg\bigr)
                                                               \tag{FF.26.10}
\]
vanishes at \(x\), by (FF.26.8)--(FF.26.9), and
its pole divisor on \(C\) has degree at most \(B\).
It must vanish identically on \(C\). All its terms
belong to the chosen finite-dimensional space, with
coefficients in the base field: here it matters that
\(\zeta_x,\eta_x\in K_p\). Injectivity of restriction
therefore gives \(R=0\) in \(K(X)\).

There is a common root \(z\) in an algebraic closure
of \(K(X)\), satisfying
\(z^d=\zeta_xh\) and \((z-1)^e=\eta_xg\).
It is nonconstant because \(h\) is nonconstant.
Every conjugate \(z'\) has the form
\(z'=\alpha z\), \(z'-1=\beta(z-1)\), for
roots of unity \(\alpha,\beta\). Consequently
\((\alpha-\beta)z=1-\beta\). As \(z\) is
transcendental over the algebraic closure of the
constant field, this forces \(\alpha=\beta=1\).
Characteristic zero makes the algebraic extension
separable, so \(z\in K(X)\). Its power is a unit
on the normal scheme \(X\), hence \(z\) is a
unit on \(X\).

The geometric group \(H^1(X_{\bar K},\mathbb Z_p(1))\)
is torsion-free. From \(d\lambda=d\kappa_p(z)\)
we obtain \(\lambda=\kappa_p(z)\) geometrically.
Normalize \(z\) to have value \(b_0\) at \(x_0\).
Hochschild--Serre then yields a unit, now denoted \(h\),
with
\[
                             \xi=\kappa_p(h).            \tag{FF.26.11}
\]
This conclusion is proved whenever the geometric class
is nonzero, without assuming in advance that the
class of \(T-1\) is nonzero. The normalized unit
also descends through every finite extension made in
this construction, before its actual point function
has been identified. For a Galois conjugate, the ratio
of the two units has zero geometric \(p\)-Kummer class,
so every coefficient of its divisor on \(\bar X\)
is divisible by every \(p^n\), and hence zero.
Properness makes this ratio constant; its value at
\(x_0\) is one. Galois descent and then evaluation
at \(x_0\) in Hochschild--Serre establish
(FF.26.11) over the original base field.

\emph{5. The zero geometric class forces a constant point function.}
If \(\lambda=0\), normalization gives
\(\xi=\kappa_p(b_0)\). By (FF.24.1), every
actual value is in \(b_0\mu_{p'}\subset K\),
including values at points over arbitrary finite
extensions. Choose a proper finite extension
\(K'/K\) and \(\beta\in K'\setminus K\).
Since no value equals \(\beta\), the pullback of
\(T-\beta\) is defined on all of \(X_{K'}\).
If its geometric class were nonzero, Step~4 would
give a nonconstant unit \(g\) with
\(b(x)-\beta=\eta_xg(x)\),
\(\eta_x\in\mu_{p'}\subset K\). All values of
\(g\) would then lie in \(K'\), contradicting
nonconstancy by the large-degree-point argument.
Thus this geometric class is zero too. Normalization
at \(x_0\) gives
\[
              \frac{b(x)}{b_0}\in\mu_{p'},\qquad
              \frac{b(x)-\beta}{b_0-\beta}\in\mu_{p'}.
\]
The two scalar coefficients belong to \(K\),
and \(1,\beta\) are linearly independent over
\(K\). These equations force \(b(x)=b_0\) at
every geometric point.

\emph{6. Removing the remaining pointwise scalars.}
We now have a unit \(h\), normalized by
\(h(x_0)=b_0\), realizing \(\xi\) whenever
the point function is nonconstant, and a constant
function otherwise. On a smooth curve with
nonconstant \(h\), (FF.24.1) gives
\(b(x)=c_xh(x)\), \(c_x\in\mu_{p'}\subset K\).
The point function of \(T-1\) on its inverse-image
open is nonconstant: otherwise the values of the
nonconstant unit \(h\) on that open would all be
in the base field, since \(c_x\in\mu_{p'}\subset K\).
The same argument applies to every \(T-a\). Thus its
class is nonzero by
Step~5. It has a unit representative \(h_1\)
and \(b(x)-1=d_xh_1(x)\), with \(d_x\in\mu_{p'}\).
Choose a point of degree greater than both a common
pole-degree bound for \(c_xh-d_xh_1-1\) and the
degree of \(h:\bar C\to\mathbb P^1\), so that
\(h(x)\notin K\). The
same divisor argument gives an identity
\[
                           h_1=Ah+B,
                              \qquad A,B\in K^\times.
\]
At every point with \(h(x)\notin K\), comparison
of the coefficients of \(1,h(x)\) forces
\(c_x=c=-A/B\).

For each \(a\in K^\times\), repeat the argument
with \(T-a\). On points with \(h(x)\notin K\),
the known value \(b(x)=ch(x)\) shows that its
unit representative is proportional to \(ch-a\).
To justify the identity rather than merely pointwise
equality, choose a point of degree above both the common
pole bound and the degree of \(h\), hence with
\(h(x)\notin K\); it exists because
\(h\) is nonconstant. At a point with
\(h(x)\in K\), set \(a=ch(x)\). If
\(b(x)\ne a\), this point belongs to the open
where \(T-a\) pulls back to a unit, but that
unit is proportional to \(ch-a\) and vanishes
there, a contradiction. Hence \(b=ch\) at
every point. Normalization at \(x_0\) gives
\(c=1\).

For higher-dimensional \(X\), pass to a finite
extension and join \(x_0\) to an arbitrary geometric
point by the normalization of a geometrically
irreducible curve through them, using
\citep[Corollary~1.9]{FF_CharlesPoonen2016} again.
The two points lift after another finite extension.
The restricted class is the Kummer class of
\(h|_C\). If this is constant, Step~5 applies;
if nonconstant, the curve argument applies. Its
normalized unit agrees with \(h|_C\): a ratio
with zero \(p\)-Kummer class has every divisor
coefficient divisible by all \(p^n\), hence is
constant, and its value at the lifted \(x_0\)
is one. Thus \(b(x)=h(x)\) at every geometric
point of \(X\).

\emph{7. Reconstruction and descent.}
The finite-extension descent in the first paragraph
puts the resulting unit over the original field.
We have proved exactly the point-function condition
for every regular connected affine source with a
rational point in
\citet[Theorem~4.18(5)]{FF_CHW2024}.
That theorem gives \(\mathcal R(K)\) for all
finite-type schemes and all pinned relative topos
morphisms, with absolute weak normalization as
in the question. Finally \(K\) is algebraic over
each of its subfields, so
Claim~\ref{clm:FFdescent} gives the asserted
subfield consequence.
\end{proof}
